% ============================================================
% harryopo-tikz-diagram - 分层架构图模板 v2 (Layered Architecture)
% ============================================================
%
% 【用途】
%   绘制 2-5 层的分层架构图，每层包含若干模块，层之间用箭头连接。
%   适用于系统架构、技术栈分层、网络协议栈等场景。
%
% 【设计原则】（参考 example-paper-structure 布局经验）
%   1. 无阴影设计：干净扁平，避免视觉噪音
%   2. 居中对齐：模块内容居中，层次清晰
%   3. 宽松内边距：inner sep=8pt，不拥挤
%   4. 统一无衬线：Arial + Microsoft YaHei
%   5. 层标题浮于虚线框上方，不打断边框
%   6. 图例嵌入 tikzpicture 内部，避免额外排版间距
%
% 【使用方法】
%   将本文件内容复制到你的 LaTeX 文档中，或使用 \input 命令引入。
%   修改下方"可调整参数"部分来自定义你的架构图。
%
% 【编译要求】
%   XeLaTeX 编译器，需加载 ctex + fontspec 宏包以支持中文。
%
% ============================================================

% ------------------------------------------------------------
% 必要的宏包（如果文档中已加载可省略）
% ------------------------------------------------------------
% \usepackage{tikz}
% \usetikzlibrary{positioning, fit, backgrounds, arrows.meta, calc}

% ------------------------------------------------------------
% 可调整参数 —— 在这里修改你的架构图配置
% ------------------------------------------------------------

% 主题色定义（Nature 无衬线风格，与 paper-structure 一致）
\definecolor{PaperBorder}{HTML}{B0C4DE}      % 柔和灰蓝边框
\definecolor{PaperFill}{HTML}{FFFFFF}        % 纯白模块填充
\definecolor{PaperTitle}{HTML}{2C3E50}       % 深蓝灰标题
\definecolor{PaperLayerBg}{HTML}{F5F8FA}     % 层背景：极淡蓝灰
\definecolor{PaperLayerBorder}{HTML}{D8E3EE} % 层边框：更浅的灰蓝
\definecolor{PaperMuted}{HTML}{7F8C8D}       % 次要文字：柔和灰
\definecolor{AccentOrange}{HTML}{D35400}     % 代码/技术栈强调：柔和橙

% 全局间距控制
\def\archLayerGap{1.6cm}      % 层与层之间的垂直距离
\def\archModuleGap{0.55cm}    % 同一层内模块之间的水平距离
\def\archLayerPad{12pt}       % 层背景框内边距（x/y统一）
\def\archTitleOffset{4pt}     % 层标题相对于背景框顶部的偏移

% 模块尺寸
\def\archModuleW{4.0cm}       % 普通模块的宽度（text width）
\def\archModuleH{1.5cm}       % 普通模块的最小高度（minimum height）
\def\archWideW{13.0cm}        % 宽模块的宽度（横跨整层）
\def\archWideH{0.95cm}        % 宽模块的最小高度

% 箭头样式
\def\archArrowThickness{1.0pt}

% ------------------------------------------------------------
% 样式定义（一般不需要修改）
% ------------------------------------------------------------
\tikzset{
  % ===== 通用模块样式 =====
  arch-module/.style = {
    rectangle,
    draw = PaperBorder,
    fill = PaperFill,
    line width = 0.8pt,
    rounded corners = 5pt,
    text width = \archModuleW,
    minimum height = \archModuleH,
    align = center,
    font = \small\sffamily,
    text = PaperTitle,
    inner sep = 8pt,
  },
  % ===== 模块内容子样式（在 node 内用 {} 包裹使用） =====
  mod-title/.style = {
    font = \small\sffamily\bfseries,
    text = PaperTitle,
  },
  mod-desc/.style = {
    font = \footnotesize\sffamily,
    text = PaperMuted,
  },
  mod-tech/.style = {
    font = \footnotesize\sffamily\ttfamily,
    text = AccentOrange,
  },
  % ===== 宽模块样式（横跨整层） =====
  arch-wide-module/.style = {
    arch-module,
    text width = \archWideW,
    minimum height = \archWideH,
  },
  % ===== 层背景框样式 =====
  arch-layerbox/.style = {
    draw = PaperLayerBorder,
    fill = PaperLayerBg,
    line width = 0.7pt,
    dashed,
    dash pattern = on 4pt off 3pt,
    rounded corners = 8pt,
    inner sep = \archLayerPad,
  },
  % ===== 层标题样式（浮于框上方，不 fill=white 打断边框） =====
  arch-layertitle/.style = {
    font = \normalsize\sffamily\bfseries,
    text = PaperTitle,
    anchor = north west,
    inner sep = 0pt,
    outer sep = 0pt,
  },
  arch-layertag/.style = {
    font = \footnotesize\sffamily\itshape,
    text = PaperMuted,
    anchor = north east,
    inner sep = 0pt,
    outer sep = 0pt,
  },
  % ===== 箭头样式（双向/下行/上行） =====
  arch-arrow/.style = {
    Stealth-Stealth,
    draw = PaperBorder,
    line width = \archArrowThickness,
  },
  arch-arrow-down/.style = {
    -Stealth,
    draw = PaperBorder,
    line width = \archArrowThickness,
  },
  arch-arrow-up/.style = {
    -Stealth,
    draw = PaperMuted,
    line width = \archArrowThickness,
  },
  % ===== 箭头标签 =====
  arr-label/.style = {
    font = \footnotesize\sffamily\bfseries,
    fill = white,
    text = PaperBorder,
    inner sep = 2pt,
    outer sep = 0pt,
  },
  arr-label-rev/.style = {
    font = \footnotesize\sffamily\bfseries,
    fill = white,
    text = PaperMuted,
    inner sep = 2pt,
    outer sep = 0pt,
  },
  % ===== 图例框 =====
  arch-legend/.style = {
    draw = PaperLayerBorder,
    fill = PaperFill,
    line width = 0.6pt,
    rounded corners = 5pt,
    inner sep = 8pt,
    font = \footnotesize\sffamily,
    text = PaperTitle,
  },
}

% ------------------------------------------------------------
% 模板正文 —— tikzpicture 环境
% ------------------------------------------------------------
%
% 【使用说明】
%   1. 第一层（最上层）的第一个节点用绝对坐标 at (0,0)
%   2. 同一层内的其他节点用 right=\archModuleGap of <节点名> 相对定位
%   3. 下一层的第一个节点用 below=\archLayerGap of <上一层>.south west, anchor=north west 定位
%   4. 每层结束后用 fit 库创建背景框和标题
%   5. 最后绘制层之间的箭头
%   6. 模块内容用 {\mod-title 标题}\\[2pt]{\mod-desc 描述}\\[1pt]{\mod-tech 技术栈} 结构
%
% ------------------------------------------------------------

\begin{tikzpicture}[
    node distance = 0pt and \archModuleGap,
    >=Stealth,
    line width = 0.7pt,
    every node/.style = {inner sep = 0pt, outer sep = 0pt},
  ]

  % ==========================================================
  % 第 1 层 —— 表示层（示例）
  % ==========================================================
  \node[arch-module] (ui-web) at (0,0) {
    {\mod-title Web 前端}
  };
  \node[arch-module, right=\archModuleGap of ui-web] (ui-mobile) {
    {\mod-title 移动 App}
  };
  \node[arch-module, right=\archModuleGap of ui-mobile] (ui-admin) {
    {\mod-title 管理后台}
  };

  \begin{scope}[on background layer]
    \node[arch-layerbox, fit=(ui-web)(ui-mobile)(ui-admin)] (layer1-box) {};
  \end{scope}
  \node[arch-layertitle] at ($(layer1-box.north west)+(0,\archTitleOffset)$) {表示层};

  % ==========================================================
  % 第 2 层 —— 业务逻辑层（示例，含多行内容）
  % ==========================================================
  \node[arch-module, below=\archLayerGap of ui-web.south west, anchor=north west] (biz-user) {
    {\mod-title 用户服务}\\[2pt]
    {\mod-desc 认证、授权}\\[1pt]
    {\mod-desc 权限管理}
  };
  \node[arch-module, right=\archModuleGap of biz-user] (biz-order) {
    {\mod-title 订单服务}\\[2pt]
    {\mod-desc 订单创建}\\[1pt]
    {\mod-tech 状态机}
  };
  \node[arch-module, right=\archModuleGap of biz-order] (biz-pay) {
    {\mod-title 支付服务}\\[2pt]
    {\mod-tech 支付宝/微信}\\[1pt]
    {\mod-desc 交易处理}
  };

  \begin{scope}[on background layer]
    \node[arch-layerbox, fit=(biz-user)(biz-order)(biz-pay)] (layer2-box) {};
  \end{scope}
  \node[arch-layertitle] at ($(layer2-box.north west)+(0,\archTitleOffset)$) {业务逻辑层};

  % ==========================================================
  % 第 3 层 —— 数据访问层（示例，含宽模块）
  % ==========================================================
  \node[arch-wide-module, below=\archLayerGap of biz-order.south, anchor=north] (data-db) {
    {\mod-title MySQL 数据库}
  };

  \begin{scope}[on background layer]
    \node[arch-layerbox, fit=(data-db)] (layer3-box) {};
  \end{scope}
  \node[arch-layertitle] at ($(layer3-box.north west)+(0,\archTitleOffset)$) {数据访问层};

  % ==========================================================
  % 层间箭头连接（示例：下行箭头 + 标签）
  % ==========================================================
  \draw[arch-arrow-down] (ui-web.south) -- (biz-user.north);
  \draw[arch-arrow-down] (ui-mobile.south) -- (biz-order.north);
  \draw[arch-arrow-down] (ui-admin.south) -- (biz-pay.north);

  \draw[arch-arrow-down] (biz-user.south) -- (data-db.north -| biz-user.south);
  \draw[arch-arrow-down] (biz-order.south) -- (data-db.north);
  \draw[arch-arrow-down] (biz-pay.south) -- (data-db.north -| biz-pay.south);

  % ==========================================================
  % 图例（可选，放在右下角）
  % ==========================================================
  % \node[arch-legend, anchor=north east] at ($(layer3-box.south east)+(0,-0.8cm)$) {
  %   \begin{tabular}{@{}c@{\hspace{6pt}}l@{\hspace{14pt}}c@{\hspace{6pt}}l@{}}
  %     \textcolor{PaperBorder}{\rule[1pt]{14pt}{1.2pt}$>$} & {\sffamily\bfseries\small 请求} &
  %     \textcolor{PaperMuted}{\rule[1pt]{14pt}{1.2pt}$>$} & {\sffamily\bfseries\small 响应}
  %   \end{tabular}
  % };

\end{tikzpicture}

% ============================================================
% 【常见修改示例】
% ============================================================
%
% 1. 增加第 4 层：
%    复制一层的代码块，修改节点名称和内容，
%    below=\archLayerGap of <上一层第一个节点>.south west, anchor=north west。
%
% 2. 右侧英文标签：
%    在 layer-box fit 之后加：
%    \node[arch-layertag] at ($(box.north east)+(0,\archTitleOffset)$) {English Tag};
%
% 3. 带标签的层间箭头：
%    \draw[arch-arrow-down] (a.south) -- (b.north)
%      node[midway, arr-label, anchor=west, xshift=3pt] {数据流向};
%    \draw[arch-arrow-up] (c.north) -- (d.south)
%      node[midway, arr-label-rev, anchor=east, xshift=-3pt] {反向流};
%
% 4. 双向箭头：
%    \draw[arch-arrow] (node1.south) -- (node2.north);
%
% 5. 折线连接（|- 和 -|）：
%    \draw[arch-arrow-down] (nodeA.south) |- (nodeB.west);
%    \draw[arch-arrow-down] (nodeA.east) -| (nodeB.north);
%
% 6. 模块内容格式（3行推荐）：
%    {\mod-title 模块名}\\[2pt]
%    {\mod-desc 功能描述行1}\\[1pt]
%    {\mod-tech 技术栈}
%    或 {\mod-desc 功能描述行2} 视内容选择。
%
% ============================================================
